\documentclass[reqno]{amsart} 
\usepackage{amssymb,latexsym,amsmath,amscd,graphicx,setspace,amsthm,verbatim}
\usepackage[margin = 3 cm]{geometry}


\theoremstyle{plain}
\newtheorem{theorem}{Theorem}[section]
\newtheorem{proposition}{Proposition}
\newtheorem{corollary}{Corollary}
\newtheorem{lemma}{Lemma}
\newtheorem{conjecture}{Conjecture}
\newtheorem{question}{Question}
\newtheorem{problem}{Problem}
      
\theoremstyle{definition}
\newtheorem{definition}{Definition}

\newenvironment{solution}{\paragraph{\emph{Solution}.}}{\hfill$\square$}


\newenvironment{solution1}{\paragraph{\emph{Solution $1$}.}}{\hfill$\square$}
\newenvironment{solution2}{\paragraph{\emph{Solution $2$}.}}{\hfill$\square$}
\newenvironment{solution3}{\paragraph{\emph{Solution $3$}.}}{\hfill$\square$}

\begin{document} 

\title[Homework 6]{Homework 6}

\date{\today} 
\maketitle 

The goal of the first problem is to make sure we all understand the chain rule from Calculus III that I recalled in class.
\begin{problem}
Let $f:\mathbb{R}^{2} \rightarrow \mathbb{R}^{3}$ be defined via
$$(x,y) \mapsto f(x,y) = (x+y,xy,x^{2}+y^{2}), $$
and let $g:\mathbb{R}^{3} \rightarrow \mathbb{R}^{2}$ be defined via
$$(u,v,w) \mapsto g(u,v,w) = (u^{2} + v, w + uv). $$
\begin{enumerate}
\item Find $Jf(x,y)$ and $Jg(u,v,w)$.
\item Let $p=(1,2)$ and $q = f(p)$.  Write down
$$Jf(p), Jg(q), Df_{p}, \text{ and } Dg_{q}. $$
(Make sure to distinguish between a matrix and a linear transformation which is a function...)
\item Calculate
$$Jg(q) \cdot Jf(p) \text{ and } Dg_{q} \circ Df_{p}.$$
\item Find a formula for $g \circ f:\mathbb{R}^{2} \rightarrow \mathbb{R}^{2}$.
\item Calculate $J (g\circ f) (x,y)$, $J (g \circ f) (p)$ and also $D(g\circ f)_{p}$.
\item Comparing (3) and (5), do we have the equalities
$$J(g \circ f)(p) = Jg(q) \cdot Jf(p) \text{ and } D(g \circ f)_{p} = Dg_{q} \circ Df_{p}??$$
If yes, did you expect to have such equalities?  Why??
\end{enumerate}
\end{problem}
\begin{solution1}

\end{solution1}


The next problem will introduce a new surface.  It will be obtained by rotating a rotating object....(think hard about this geometrically as you go over the problem).
\begin{problem}
Remember that a counter-clockwise rotation of angle $\varphi$ in $\mathbb{R}^{2}$ can be described using the square matrix
\begin{equation*}
\begin{pmatrix}
\cos \varphi & - \sin \varphi \\
\sin \varphi & \cos \varphi
\end{pmatrix}
\end{equation*}
\begin{enumerate}
\item Describe in words what is the object parametrized by the function $\gamma:(-1,1) \rightarrow \mathbb{R}^{2}$ given by 
$$t \mapsto (t+3,t).$$  
You can even draw a picture.
\item Describe in words what is the object parametrized by the function $(-1,1) \times (0,\pi/2) \rightarrow \mathbb{R}^{2}$ given by 
$$(t, \varphi)  \mapsto (3 + t \cos(\varphi) - t\sin(\varphi), t \sin(\varphi) + t \cos(\varphi)).$$  
You can even draw a picture.
\item You should now have the image of a line segment spinning around the point $(3,0) \in \mathbb{R}^{2}$.  Imagine now that as the line segment is spinning around that point, the line segment itself is rotating around the $y$-axis just like what was happening for a surface of revolution.  Try to describe in words what is the object parametrized by the function $(-1,1) \times (0,2\pi) \mapsto \mathbb{R}^{3}$ given by
$$(t,\theta) \mapsto \left(\left(3 + t \cos\frac{\theta}{2} - t \sin \frac{\theta}{2}\right)\cos \theta,t \sin \frac{\theta}{2} + t \cos \frac{\theta}{2},\left(3 + t \cos\frac{\theta}{2} - t \sin \frac{\theta}{2}\right)\sin \theta\right). $$
Once you are done imagining with your own head, I would encourage you to go to Geogebra, google how to use the command Surface, and enter that equation in the 3D calculator to see this surface.
\item Write down two local parametrizations carefully that would give the structure of a regular $2$-dimensional smooth manifold on this surface.
\item If you are brave, go ahead pick one of the two and check that it satisfies the three usual conditions.
\end{enumerate}
\end{problem}
\begin{solution1}

\end{solution1}


















\end{document} 


